% \section{Analysis} \label{sec:analysis}

\section{Results and Analysis}

% SoM discussion
\begin{figure*}[t]
    \centering
    \includegraphics[width=\linewidth]{figures/execution_trace_v2.pdf}
    % \vspace{-0.32in}
    \caption{Successful execution trajectory of the GPT-4V + SoM agent on the task for blocking a user that posted a certain picture. The text in red represents the actions output by the agent.} \label{fig:block-user-traj}
    % \vspace{-0.1in}
\end{figure*}



Our main baseline results are summarized in Tab.~\ref{tab:results}. All existing models substantially underperform compared to humans, which indicate significant headroom in \BenchmarkName{} for future work. We discuss some main findings below with the GPT-4V model, with further analysis in the appendix.

\paragraph{Text-based LLMs Perform Poorly} State-of-the-art text-only LLMs generally achieve poor results, with the best model (GPT-4) achieving an overall success rate of 7.25\%. When we augment the LLMs with captions, this considerably improves success rate (7.25\% to 12.75\% for GPT-4). % Other models also see a similar boost in success rate when provided caption information.

\paragraph{Multimodality Helps} Using multimodal agents significantly improves the success rate: GPT-4V (\texttt{gpt-4-1106-vision-preview}) achieves an overall success rate of 15.05\%, substantially improving over the text-only agents. Gemini-Pro also experiences a significant uplift in success rate, from 3.85\% (caption-augmented) to 6.04\% (multimodal). 
% These results are intuitive, as \BenchmarkName{} was designed to be explicitly visual, and test the ability of models to process and understand visual content. 
Text-based agents may be limited in their ability to process complex images (e.g., those that require OCR or recognition of non-salient objects). %, thus falling behind multimodal models.



\paragraph{SoM Improves Navigability} We observe that the SoM representation (Sec.~\ref{sec:som_method}) further improves the performance of GPT-4V over using the accessibility tree, boosting overall success rate ($15.05\% \rightarrow 16.37\%$). We observe particularly substantial improvements on Classifieds and Reddit, from $12.38\% \rightarrow 17.14\%$ and $8.12\% \rightarrow 9.83\%$ respectively. 
% In contrast, the performance on the Shopping benchmark remains relatively similar. 
We attribute this to the Classifieds and Reddit websites containing denser visual content. These websites often contain multiple smaller sized images that are arranged very closely (Fig.~\ref{fig:som_figure}). In many of these pages, the accessibility tree does not provide sufficient information to disentangle elements that are spatially close. 
We hypothesize that the SoM representation is superior for strong VLM agents, which can more accurately disentangle and click on the desired elements. 
For the other VLMs, SoM does not significantly improve success, which we attribute to the finding from \citet{yang2023som} that only GPT-4V demonstrates the SoM grounding ability (perhaps due to scale or training data). % This motivates future work in imbueing VLM agents with similar abilities.

One GPT-4V + SoM execution trajectory that we found particularly compelling was Reddit task \#139, which requires exact image matching to find a post and block a user (Fig.~\ref{fig:block-user-traj}). The model initially attempts to search for the correct forum, and when this fails it navigates to the list of forums. After navigating correctly to \texttt{/f/memes}, it identifies the offending image out of the many images on the page (Step 3 in Fig.~\ref{fig:block-user-traj}) and blocks the author successfully without any unnecessary actions. 

\subsection{Performance by Task Type}  \label{sec:performance_task_subset}

We analyze the success rate of the best VLM agent baseline (GPT-4V + SoM) across several additional subsets of tasks (Tab.~\ref{tab:analysis_subsets}). We include further analysis for other models in Appendix~\ref{sec:appendix_analysis}.
%and discussed in the following paragraphs, with

\begin{table}[t] %{r}{6cm}
% \captionsetup{font=small}
\centering
\resizebox{0.85\linewidth}{!}{%
\begin{tabular}{lcc}
  \toprule
  \textbf{Task Subset}\hspace{25mm} & \textbf{\% of Total} & \textbf{SR} ($\uparrow$) \\
  \midrule
  OCR required &  17.1\% &  13.4\% \\
  No OCR required &  82.9\% & \textbf{16.9\%} \\
  \midrule
  Exact image match &  8.7\% &  \textbf{18.9\%} \\
  No exact image match &  91.3\% & 16.2\% \\
  \midrule
  Image inputs &  25.2\% &  \textbf{19.0\%} \\
  No image inputs &  74.8\% & 14.9\% \\
  \bottomrule
\end{tabular}
}
% \vspace{-0.1in}
\caption{Success rate (SR) of GPT-4V (SoM) across different types of tasks.}  \label{tab:analysis_subsets}
% \vspace{-0.2in}
\end{table}

\paragraph{OCR Tasks} 17.1\% of \BenchmarkName{} require optical character recognition (OCR), such as reading text from product images, or extracting text from an input image. We find that GPT-4V + SoM generally performs worse on tasks that require OCR (13.4\%) compared to tasks which do not (16.9\%), suggesting that OCR may be a bottleneck for current agents.

\paragraph{Exact Image Match} 8.7\% of tasks require exact image matching, which requires agents to identify precise visual matches. GPT-4V + SoM achieves a slightly higher success rate on this subset (18.9\%) compared to other tasks (16.2\%), suggesting that exact image matching is not a primary bottleneck.


\paragraph{Image Input Tasks} 25.2\% of \BenchmarkName{} include one or more input images as part of the objective. These tasks generally appear more tractable for the GPT-4V + SoM agent, and it achieves a higher success rate (19.0\%) compared to tasks without image inputs (14.9\%).% In contrast, Gemini-Pro performance deteriorates, and it performs 



